% Auto-generated from meeting_2569-05-16_semester3-closeout-and-satisfaction.md (กบ.วช. format v2)
\documentclass{meetingminutes}
\meetingcommittee{คณะกรรมการบริหารหลักสูตร วิศวกรรมคอมพิวเตอร์และปัญญาประดิษฐ์ (บัณฑิตศึกษา)}
\meetingnumber{๒/ภาคเรียนที่ ๓/๒๕๖๘}{๒๕๖๘}
\meetingdate{วันที่ ๑๖ พฤษภาคม พ.ศ. ๒๕๖๙}
\meetingplace{ห้องประชุมคณะเทคโนโลยีอุตสาหกรรม}
\meetingstart{๑๓.๐๐ น.}
\meetingend{๑๕.๔๕ น.}

\begin{document}
\maketitle

\psection{ผู้มาประชุม}
\begin{participants}
  \pmember{1}{ผู้ช่วยศาสตราจารย์ ดร.พัชรี มณีรัตน์}{ประธานหลักสูตร}{ประธาน}
  \pmember{2}{ผู้ช่วยศาสตราจารย์ ดร.พิศิษฐ์ นาคใจ}{อาจารย์ประจำหลักสูตร}{กรรมการ}
  \pmember{3}{ผู้ช่วยศาสตราจารย์ ดร.กาญจนา ดาวเด่น}{อาจารย์ประจำหลักสูตร (เลขานุการ)}{กรรมการ/เลขานุการ}
\end{participants}

\psection{ผู้ไม่มาประชุม}
\noindent ไม่มี\par

\startmeeting
\openingchair{ผู้ช่วยศาสตราจารย์ ดร.พัชรี มณีรัตน์}{กล่าวเปิดการประชุมและดำเนินการประชุมตามระเบียบวาระ ดังนี้}

\agenda{๑}{รับรองรายงานการประชุมครั้งที่ 1/ภาค 3/2568}
ที่ประชุมรับรองรายงานการประชุมครั้งที่ 1/ภาค 3/2568 (25 เมษายน 2569) โดยไม่มีการแก้ไข


\agenda{๒}{สรุปผลการดำเนินงานหลักสูตรปีการศึกษา 2568 (ภาพรวม)}
ประธานนำเสนอสรุปผลการดำเนินงานหลักสูตรในปีการศึกษา 2568:

ด้านการรับนักศึกษา: - รุ่นที่ 1: สมัคร 6 คน รายงานตัว 4 คน (66.7\%) - นักศึกษาปีการศึกษา 2569: สมัคร 8 คน (เพิ่มขึ้น 33\%)

ด้านผลการเรียน:

\begin{itemize}
  \item ภาค ๒/๒๕๖๘ (ภาคเรียนที่ ๑ ของหลักสูตร — รายวิชา ๔๐๙๕๑๐๑ และ ๔๐๙๕๑๐๒): นักศึกษาผ่านทุกรายวิชา Pass Rate = ๑๐๐ เปอร์เซ็นต์
  \item ภาค ๓/๒๕๖๘ (ภาคเรียนที่ ๒ ของหลักสูตร — รายวิชา ๗๐๑๕๑๐๑ และ ๗๐๑๕๙๐๖): นักศึกษาผ่านทุกรายวิชา Pass Rate = ๑๐๐ เปอร์เซ็นต์
  \item ไม่มีนักศึกษาลาออกหรือพักการศึกษาตลอดปีการศึกษา ๒๕๖๘
\end{itemize}

ด้านผลงาน/ความสำเร็จ: - นายจิรกิตติ์ วราภรณ์: รางวัลชนะเลิศ + ทองแดงภาคเหนือ (Remote Sensing) - ผศ.ดร.พิศิษฐ์: ทุน Fundamental Fund + ทุน สวก. - กิจกรรม IEL Community Workshop: 3 โรงเรียน 26 มี.ค. 2569


\agenda{๓}{นำเสนอผลการสำรวจความพึงพอใจนักศึกษา (AUNQA-C6-STSAT-2568)}
ผศ.ดร.พัชรี มณีรัตน์ นำเสนอผลการสำรวจความพึงพอใจนักศึกษาปัจจุบัน (n=4):

\begin{longtable}{|p{6cm}|p{3cm}|p{3cm}|}
  \hline
  \textbf{ด้าน} & \textbf{คะแนนเฉลี่ย} & \textbf{ระดับ} \\
  \hline
  ด้านเนื้อหาและหลักสูตร (C1/C2) & 4.97/5.00 & ดีเยี่ยม \\
  \hline
  ด้านการจัดการเรียนการสอน (C3) & 4.93/5.00 & ดีเยี่ยม \\
  \hline
  ด้านการประเมินผล (C4) & 4.95/5.00 & ดีเยี่ยม \\
  \hline
  ด้านอาจารย์ผู้สอน (C5) & 4.97/5.00 & ดีเยี่ยม \\
  \hline
  ด้านบริการสนับสนุน (C6) & 4.90/5.00 & ดีเยี่ยม \\
  \hline
  ด้านสิ่งอำนวยความสะดวก (C7) & 4.93/5.00 & ดีเยี่ยม \\
  \hline
  คะแนนเฉลี่ยรวม & 4.94/5.00 & ดีเยี่ยม \\
  \hline
\end{longtable}

100\% ของนักศึกษาเห็นว่าหลักสูตรสอดคล้องกับตลาดแรงงาน

ที่ประชุมรับทราบผลการสำรวจและเห็นว่าเป็นหลักฐานสำคัญสำหรับ AUN-QA R8.5


\agenda{๔}{ทบทวนและวางแผนพัฒนาหลักสูตรต่อไป}
ที่ประชุมอภิปรายแผนพัฒนาหลักสูตรในปีการศึกษา 2569: 1. เพิ่มความร่วมมือกับภาคอุตสาหกรรม (MOU/Industry Project) 2. ออกแบบแบบสอบถามความพึงพอใจสำหรับผู้ใช้บัณฑิต 3. เพิ่มวิชาเลือกด้าน Generative AI และ LLM 4. ส่งเสริมการตีพิมพ์บทความวิจัยนักศึกษา (เป้าหมาย: TCI หรือ Scopus/WoS ตามเกณฑ์ กมอ.)

\resolution{รับทราบทุกวาระ และอนุมัติแผนพัฒนาหลักสูตรปีการศึกษา 2569 ตามที่เสนอ}

\adjournmeeting

\begin{sigblock}
\typist{ผู้ช่วยศาสตราจารย์ ดร.กาญจนา ดาวเด่น}
\reviewer{ผู้ช่วยศาสตราจารย์ ดร.พัชรี มณีรัตน์}{กรรมการ}
\chairsig{ผู้ช่วยศาสตราจารย์ ดร.พัชรี มณีรัตน์}{ประธานการประชุม}
\end{sigblock}

\end{document}
